\documentclass[11pt,a4paper,onecolumn,final]{article}

%OT1 is the default font encoding of memoir class
%It is generally advised to put it before inputenc
%\usepackage{kerkis}
\usepackage[T1]{fontenc}
\usepackage{graphicx}
\usepackage{float}
\usepackage{chngcntr}
\counterwithin{figure}{section}

\usepackage{tikz}
\usetikzlibrary{angles, quotes}

\usepackage{caption}
%Input encoding for typing directly Greek characters
%Other values for the option: latin1, ascii, ansinew
%It is generally advised to put it before babel
\usepackage[utf8]{inputenc}

%MAIN Language the LAST one
\usepackage[greek,english]{babel}


\usepackage[
backend=biber,
style=alphabetic,
sorting=ynt
]{biblatex}
\addbibresource{sample.bib}

%\usepackage{kerkis}

%The teubner package extends the facilities offered by the CB Greek fonts
%NOTE: This includes also the graphicx package!
%NOTE: When the teubner package is loaded, command \cap (denoting set intersection) is overwritten
%NOTE: TEUBNER.STY FILE HAS BEEN CHANGED A BIT IN LINES 255-257 WITH 3 REPLACEMENTS \cap -> \capt
\usepackage{teubner}

%Package for displaying math
\usepackage{amsmath}

%Loading a superset of amsfonts
\usepackage{amssymb}

%Provides an enhanced version of \newtheorem
%1.Defines * for unnumbered environments
%2.Defines a proof environment
%3.Defines plain, definition, and remark theorem styles
\usepackage{amsthm}

%Allows text in eps files to be replaced with LATEX symbols
\usepackage{psfrag}

%An extension of the cite package, compatible with the IEEEtranN bibliography style
\usepackage[numbers,sort&compress]{natbib}

\makeatletter\let\c@lofdepth\relax\let\c@lotdepth\relax\makeatother
\usepackage[small,bf,tight]{subfigure}

%Math symbols of the dsfont style. Enables the use of \mathds{}
\usepackage{dsfont}

%Defines several symbols in text mode (e.g. the \textreferencemark for use in the itemize environment)
\usepackage{textcomp}

%Enables the Ralph Smith's Formal Script font in math mode via the \mathscr{} command
\usepackage{mathrsfs}

\usepackage{geometry}
%PDF VIEW
\geometry{total={210mm,297mm},left=25mm,right=25mm,bindingoffset=0mm, top=25mm,bottom=25mm}
%PRINT
%\geometry{total={210mm,297mm},left=20mm,right=20mm,bindingoffset=10mm, top=25mm,bottom=25mm}

\usepackage[breaklinks=true,colorlinks=true,
%linkcolor=blue,urlcolor=blue,citecolor=blue,% PDF VIEW
linkcolor=black,urlcolor=black,citecolor=black,% PRINT
bookmarks=false,bookmarksopenlevel=2]{hyperref}
\usepackage{hyperref}

\usepackage{cleveref}

%Convenient language switch
\newcommand{\gr}[1]{\foreignlanguage{greek}{#1}}
\newcommand{\en}[1]{\foreignlanguage{english}{#1}}

%\newtheorem{definition}{Ορισμός}
\newtheorem{definition}{Definition}
\newtheorem{example}{Example}
\newtheorem{exercise}{Exercise}
%\newtheorem{lemma}{Λήμμα}
\newtheorem{lemma}{Lemma}
%\newtheorem{proof}{Proof}
\newtheorem{notation}{Notation}
\newtheorem{problem}{Problem}
%\newtheorem{proposition}{Πρόταση}
\newtheorem{proposition}{Proposition}
\newtheorem{remark}{Remark}
\newtheorem{solution}{Solution}
\newtheorem{summary}{Summary}
\newcommand{\argmax}{\arg\!\max}
%\addto\captionsgreek{\renewcommand{\chaptername}{Κεφάλαιο}}
%\addto\captionsgreek{\renewcommand{\tablename}{Πίνακας }}
%\addto\captionsgreek{\renewcommand{\proofname}{Απόδειξη }}

\makeatletter
%\renewcommand{\fnum@figure}{Σχήμα \thefigure}
%\renewcommand{\fnum@table}{Πίνακας \thetable}
\makeatother
\setlength{\textfloatsep}{1pt}

\begin{document}

%COVER PAGE%%%%%%%%%%%%%%%%%%%%%%%%%%%%%%%%%%%%%%%%%%%%%%%%%%%%%%%%%%%%%%%%%%%%%
\thispagestyle{empty}
{\sffamily\centering\Large

\vspace{1.0cm}

{\Large \textbf{3η Εργαστηριακή άσκηση}\\
\textbf{Μέθοδος Μέγιστης Καθόδου με Προβολή} \\[0.5cm]

\vspace{1.0cm}

{\LARGE Λαβδάκης Γεώργιος \\
ΑΕΜ:10672}

\vspace{1.0cm}

\tableofcontents

\newpage

\begin{flushleft}

Σκοπός της 3ης εργασίας είναι η ελαχιστοποίηση της συνάρτησης \(f(x_{1},x_{2})=\frac{1}{3}x_{1}^2+3x_{2}^2\) αρχικά, όταν κανένας περιορισμός δεν εφαρμόζεται στις μεταβλητές \(x\), με τη μέθοδο της μέγιστης καθόδου και μετά, θεωρώντας τους περιορισμούς \(-10\leq x_{1} \leq 5\) και \(-8\leq x_{2} \leq 12\), με τη μέθοδο της μέγιστης καθόδου με προβολή. Αφού υλοποιείται η κάθε μέθοδος, σχολιάζονται τα σχετικά αποτελέσματα.

\section{Θέμα 1}

Σε αυτό το θέμα θα υλοποιήσουμε τη μέθοδο της μέγιστης καθόδου για να ελαχιστοποιήσουμε τη συνάρτηση \(f(x_{1},x_{2})\) με ακρίβεια \(\epsilon=0.001\) και σημείο εκκίνησης (-4,4) για διάφορες σταθερές τιμές του βήματος \(\gamma_{k}\). Αρχικά, σχεδιάζουμε την \(f\), για να πάρουμε μια γενική εικόνα της συνάρτησης, στο σχήμα \ref{fig:plot_func}.

\begin{figure}[H]
    \centering
    \includegraphics[width=0.8\linewidth]{plot_func.jpg}
    \caption{Γραφική παράσταση της \(f(x_{1},x_{2})=\frac{1}{3}x_{1}^2+3x_{2}^2\)}
    \label{fig:plot_func}
\end{figure}

Για τα επόμενα ερωτήματα θα χρειαστούμε το grad της \(f(x_{1},x_{2})\), \(\nabla f(x)=\begin{bmatrix} \frac{2x_{1}}{3} \\6x_{2}\end{bmatrix}\).

\newpage

Για σταθερό βήμα \(\gamma_{k}=0.1\) προκύπτουν οι γραφικές παραστάσεις \ref{fig:steep1_i} και \ref{fig:steep2_i}.

\begin{figure}[H]
    \centering
    \includegraphics[width=0.7\linewidth]{steep1_i.jpg}
    \caption{Εύρεση ελαχίστου με τη μέθοδο της μέγιστης καθόδου για σταθερό \(\gamma_{k}=0.1\) με σημείο εκκίνησης (-4,4)}
    \label{fig:steep1_i}
\end{figure}

\begin{figure}[H]
    \centering
    \includegraphics[width=0.8\linewidth]{steep2_i.jpg}
    \caption{Τιμή συνάρτησης κατά την εύρεση ελαχίστου συναρτήσει των επαναλήψεων}
    \label{fig:steep2_i}
\end{figure}

Όμοια, για σταθερό βήμα \(\gamma_{k}=0.3\) προκύπτουν οι γραφικές παραστάσεις \ref{fig:steep1_ii} και \ref{fig:steep2_ii}.

\begin{figure}[H]
    \centering
    \includegraphics[width=0.7\linewidth]{steep1_ii.jpg}
    \caption{Εύρεση ελαχίστου με τη μέθοδο της μέγιστης καθόδου για σταθερό \(\gamma_{k}=0.3\) με σημείο εκκίνησης (-4,4)}
    \label{fig:steep1_ii}
\end{figure}

\begin{figure}[H]
    \centering
    \includegraphics[width=0.8\linewidth]{steep2_ii.jpg}
    \caption{Τιμή συνάρτησης κατά την εύρεση ελαχίστου συναρτήσει των επαναλήψεων}
    \label{fig:steep2_ii}
\end{figure}

Με παρόμοιο τρόπο, για σταθερό βήμα \(\gamma_{k}=3\) προκύπτουν οι γραφικές παραστάσεις \ref{fig:steep1_iii} και \ref{fig:steep2_iii}.

\begin{figure}[H]
    \centering
    \includegraphics[width=0.7\linewidth]{steep1_iii.jpg}
    \caption{Εύρεση ελαχίστου με τη μέθοδο της μέγιστης καθόδου για σταθερό \(\gamma_{k}=3\) με σημείο εκκίνησης (-4,4)}
    \label{fig:steep1_iii}
\end{figure}

\begin{figure}[H]
    \centering
    \includegraphics[width=0.8\linewidth]{steep2_iii.jpg}
    \caption{Τιμή συνάρτησης κατά την εύρεση ελαχίστου συναρτήσει των επαναλήψεων}
    \label{fig:steep2_iii}
\end{figure}

Τέλος, για σταθερό βήμα \(\gamma_{k}=5\) προκύπτουν οι γραφικές παραστάσεις \ref{fig:steep1_iv} και \ref{fig:steep2_iv}.

\begin{figure}[H]
    \centering
    \includegraphics[width=0.7\linewidth]{steep1_iv.jpg}
    \caption{Εύρεση ελαχίστου με τη μέθοδο της μέγιστης καθόδου για σταθερό \(\gamma_{k}=5\) με σημείο εκκίνησης (-4,4)}
    \label{fig:steep1_iv}
\end{figure}

\begin{figure}[h]
    \centering
    \includegraphics[width=0.8\linewidth]{steep2_iv.jpg}
    \caption{Τιμή συνάρτησης κατά την εύρεση ελαχίστου συναρτήσει των επαναλήψεων}
    \label{fig:steep2_iv}
\end{figure}

\(\cdot\) Παρατηρούμε ότι για τις δύο αρχικές τιμές του βήματος \(\gamma_{k}=0.1\) και \(\gamma_{k}=0.3\) ο αλγόριθμος συγκλίνει ενώ, για τις δύο τελευταίες \(\gamma_{k}=3\) και \(\gamma_{k}=5\) όχι. Θα αποδείξουμε τα παραπάνω αποτελέσματα βρίσκοντας τις τιμές του \(\gamma_{k}\) για τις οποίες συγκλίνει ο αλγόριθμος.

\(\cdot\) Για να συγκλίνει ο αλγόριθμος πρέπει να ισχύει \(\frac{|x_{k+1}|}{|x_{k}|}<1\). Δηλαδή, πρέπει κάθε επόμενη τιμή της ακολουθίας \(\{x_{k}\}\) να είναι μικρότερη από την προηγούμενη, έτσι ώστε να συγκλίνει στο ελάχιστο, το οποίο είναι προφανώς (0,0) καθώς η αντικειμενική συνάρτηση \(f(x_{1},x_{2})=\frac{1}{3}x_{1}^2+3x_{2}^2\) είναι άθροισμα των τετραγώνων των 2 μεταβλητών. Για τη μέθοδο της μέγιστης καθόδου ισχύει \(x_{k+1}=x_{k}-\gamma_{k}\nabla f(x_{k})\) και πρέπει οι ακολουθίες για το \(x_{1}\) και για το \(x_{2}\) να συγκλίνουν ταυτόχρονα στο μηδέν. Αφού έχουμε υπολογίσει το \(\nabla f(x)\) παραπάνω, προκύπτει το σύστημα: 

\(x_{1k+1}=x_{1k}-\gamma_{k}\frac{2x_{1k}}{3} \Rightarrow \frac{|x_{1k+1}|}{|x_{1k}|}=|1-\frac{\gamma_{k}}{3}|<1 \Rightarrow 6>\gamma_{k}>0 \)

\(x_{2k+1}=x_{2k}-\gamma_{k}6x_{2k} \Rightarrow \frac{|x_{2k+1}|}{|x_{2k}|}=|1-6\gamma_{k}|<1 \Rightarrow \frac{1}{3}>\gamma_{k}>0\)

Άρα, για να συγκλίνει ο αλγόριθμος πρέπει να ισχύει \(0<\gamma_{k}<\frac{1}{3}\) γεγονός που εξηγεί γιατί συγκλίνει στις περιπτώσεις που \(\gamma_{k}=0.1<\frac{1}{3}\) και \(\gamma_{k}=0.3<\frac{1}{3}\) ενώ, δεν συγκλίνει για \(\gamma_{k}=3>\frac{1}{3}\) ή \(\gamma_{k}=5>\frac{1}{3}\).

\(\cdot\) Η υλοποίηση της μεθόδου της μέγιστης καθόδου βρίσκεται στο αρχείο steepest\_descent.m.
 
\section{Θέμα 2}

Στα επόμενα θέματα θα υλοποιήσουμε τη μέθοδο της μέγιστης καθόδου με προβολή για να ελαχιστοποιήσουμε τη συνάρτηση \(f(x_{1},x_{2})\) με περιορισμούς \(-10\leq x_{1} \leq 5\) και \(-8\leq x_{2} \leq 12\) και ακρίβεια \(\epsilon=0.01\) μεταβάλλοντας το βήμα προς την κατεύθυνση της αρνητικής κλίσης \(s_{k}\), το βήμα της μέγιστης καθόδου \(\gamma_{k}\) και το σημείο εκκίνησης. Σε αυτό το θέμα θα υλοποιήσουμε τη μέθοδο για \(\gamma_{k}=0.5, s_{k}=5\) και σημείο εκκίνησης (5,-5). Προκύπτουν οι γραφικές παραστάσεις \ref{fig:sd_proj_2i} και \ref{fig:sd_proj_2ii}.

\begin{figure}[H]
    \centering
    \includegraphics[width=0.7\linewidth]{sd_proj_2i.jpg}
    \caption{Εύρεση ελαχίστου με τη μέθοδο της μέγιστης καθόδου με προβολή για σταθερό \(\gamma_{k}=0.5\), \(s_{k}=5\) και σημείο εκκίνησης (-5,5)}
    \label{fig:sd_proj_2i}
\end{figure}

\begin{figure}[H]
    \centering
    \includegraphics[width=0.8\linewidth]{sd_proj_2ii.jpg}
    \caption{Τιμή συνάρτησης κατά την εύρεση ελαχίστου συναρτήσει των επαναλήψεων}
    \label{fig:sd_proj_2ii}
\end{figure}

Παρατηρούμε ότι το σημείο εκκίνησης είναι εφικτό σημείο και καθώς εκτελείται ο αλγόρθμος μένουμε στην περιοχή των εφικτών σημείων. Επίσης, το βήμα είναι \(s_{k} \cdot \gamma_{k}=2.5>\frac{1}{3}\) άρα, σύμφωνα με την ανάλυση που κάναμε στο θέμα 1 ο αλγόριθμος δεν συγκλίνει, όπως φαινεται και από τη \ref{fig:sd_proj_2ii} και υπάρχει ταλάντωση μεταξύ δύο σημείων. 

\newpage

\section{Θέμα 3}

Σε αυτό το θέμα θα υλοποιήσουμε τη μέθοδο για \(\gamma_{k}=0.1, s_{k}=15\) και σημείο εκκίνησης (-5,10). Προκύπτουν οι γραφικές παραστάσεις \ref{fig:sd_proj_3i} και \ref{fig:sd_proj_3ii}.

\begin{figure}[H]
    \centering
    \includegraphics[width=0.7\linewidth]{sd_proj_3i.jpg}
    \caption{Εύρεση ελαχίστου με τη μέθοδο της μέγιστης καθόδου με προβολή για σταθερό \(\gamma_{k}=0.1\), \(s_{k}=15\) και σημείο εκκίνησης (-5,10)}
    \label{fig:sd_proj_3i}
\end{figure}

\begin{figure}[H]
    \centering
    \includegraphics[width=0.8\linewidth]{sd_proj_3ii.jpg}
    \caption{Τιμή συνάρτησης κατά την εύρεση ελαχίστου συναρτήσει των επαναλήψεων}
    \label{fig:sd_proj_3ii}
\end{figure}

\(\cdot\) Παρατηρούμε ότι το σημείο εκκίνησης είναι εφικτό σημείο και καθώς εκτελείται ο αλγόρθμος μένουμε στην περιοχή των εφικτών σημείων. Επίσης, το βήμα είναι \(s_{k} \cdot \gamma_{k}=1.5>\frac{1}{3}\) όπως στο θέμα 2, όμως, ο αλγόριθμος πλησιάζει μια περιοχή κοντά στο ελάχιστο, ταλαντώνεται γύρω από αυτή και εν τέλει τερματίζει μόνος του.

\(\cdot\) Ένας πρακτικός τρόπος ώστε η μέθοδος να συγκλίνει στο ελάχιστο θα ήταν η επιλογή μικρότερου \(\gamma_{k}\) ή \(s_{k}\) ή επιλογή του \(\gamma_{k}\) όχι σταθρού αλλά, σύμφωνα με τον κανόνα armijo.

\section{Θέμα 4}

Σε αυτό το θέμα θα υλοποιήσουμε τη μέθοδο για \(\gamma_{k}=0.2, s_{k}=0.1\) και σημείο εκκίνησης (8,-10). Προκύπτουν οι γραφικές παραστάσεις \ref{fig:sd_proj_4i} και \ref{fig:sd_proj_4ii}.

\begin{figure}[H]
    \centering
    \includegraphics[width=0.7\linewidth]{sd_proj_4i.jpg}
    \caption{Εύρεση ελαχίστου με τη μέθοδο της μέγιστης καθόδου με προβολή για σταθερό \(\gamma_{k}=0.2\), \(s_{k}=0.1\) και σημείο εκκίνησης (8,-10)}
    \label{fig:sd_proj_4i}
\end{figure}

\begin{figure}[H]
    \centering
    \includegraphics[width=0.8\linewidth]{sd_proj_4ii.jpg}
    \caption{Τιμή συνάρτησης κατά την εύρεση ελαχίστου συναρτήσει των επαναλήψεων}
    \label{fig:sd_proj_4ii}
\end{figure}

\(\cdot\) Παρατηρούμε ότι παρόλο που το σημείο εκκίνησης δεν είναι εφικτό σημείο και ο αλγόριθμος δεν θα έπρεπε να συγκλίνει, η προβολή μας επαναφέρει στον χώρο των εφικτών σημείων και τελικά ο αλγόριθμος συγκλίνει. Αυτό οφείλεται και στην επιλογή μικρών τιμών για τα βήματα \(\gamma_{k}, s_{k}\). Με τις επιλογές αυτών των παραμέτρων ωστόσο, απαιτείται μεγάλος αριθμός επαναλήψεων για τη σύγκλιση του αλγορίθμου. 

\(\cdot\) Η υλοποίηση της μεθόδου της μέγιστης καθόδου με προβολή και η συνάρτηση που υπολογίζει την προβολή ενός διανύσματος z στους περιορισμούς X που ορίσαμε βρίσκονται στα αρχεία steepest\_descent\_projection.m και projection.m αντίστοιχα.

\end{flushleft}

\end{document}
